\chapter{Agents across the causal decision workflow}
\label{ch:la04-workflow}

\section{A decision contract before an agent architecture}

A banking assistant becomes useful when it helps an analyst preserve the meaning
of a decision across several transformations. A treasury question becomes a
population and an intervention; account records become measurements of that
population; a diagram becomes an identification problem; an identified
functional becomes a numerical request; and the result becomes a monetary
comparison. Each transformation can change the question without producing a
software error. Fluent explanations are particularly dangerous when they make
these changes hard to notice. The appropriate object of automation is therefore
a sequence of accountable transformations, with explicit unresolved states.

Consider a request to ``reduce the funding cost of the customers likely to
leave.'' An analyst must establish whether the intervention is a deposit-rate
offer, a service action, or a portfolio allocation; whether the unit is an
account, customer, household, or business group; whether likely departure is a
baseline score or a future outcome; and whether cost includes the offer paid to
customers who would have stayed anyway. An assistant can list these alternatives
and attach relevant contract passages. It cannot establish which alternative
the bank intends merely by selecting the most plausible sentence. A structured
interview, conducted without an LLM, is the comparison baseline. The relevant
quality measure is the accuracy and completeness of the resulting target under
expert review, together with analyst time and unresolved-question rates.

We represent the decision contract schematically as
\[
 Q=(\mathcal P,U,A,g_0,g_1,Y,\tau,\nu,\mathcal H,\mathcal A),
\]
where $\mathcal P$ is the population, $U$ the unit, $A$ the feasible action,
$g_0,g_1$ the comparison regimes, $Y$ the outcome, $\tau$ the horizon, $\nu$ the
target measure, $\mathcal H$ the available histories and $\mathcal A$ the declared
assumptions. An ordinary population mean and a propensity-overlap mean have
different $\nu$. A policy responding at day 11 and one acting at day 10 may have
different $\mathcal H$, even when they reference the same event. This tuple is a
project specification, not a theorem that every target is identified. Its
purpose is to make semantic changes observable.

The literature supplies several useful components but no generally qualified
substitute for this contract. Causal discovery assistants can separate user
decisions from formal algorithm outputs; retrieval systems can attach passages;
benchmark environments can require structured identification and estimation
answers. These roles have different failure modes and should not receive a
single undifferentiated ``agent accuracy'' score
\cite{v12_la_2606_23608,v12_la_2402_15301,v12_la_package_causalds}.

\section{Evidence retrieval and the construction of observations}

Retrieval has two distinct products: a discovered source and a supported claim.
A search result supplies the first. The second requires reading primary text,
checking its version and population, and showing why its assumptions apply to
the present query. A banking process manual may establish when a contact occurs
without establishing why customers receive it. A statistical paper may justify
an estimator conditional on exchangeability without establishing that
exchangeability holds in the bank. A source-to-claim relation must therefore
record both the supported statement and the applicability judgment.

The retrieval baseline is a curated collection reviewed by an analyst. A useful
assistant can broaden candidate discovery, locate conflicting passages and
produce an evidence table for review. Evaluation should include relevant,
irrelevant and misleading documents, missing sources, obsolete versions and
sources that contradict each other. Retrieval failure must remain distinguishable
from evidence that a proposed relation is false. The LACR results discussed in
the preceding chapter are a concrete reason not to assume that more retrieved
text monotonically improves discovery \cite{v12_la_2402_15301}.

Observation construction is equally substantive. A field called
\texttt{owner\_id} may identify an operational account holder rather than the
economic owner relevant to exposure aggregation. A link score is a measurement
about a proposed relation, not the relation itself. The observation journey must
retain candidate mappings, inclusion rules, missed records, joint dependence
and the origin of adjudicated labels. Ontology tools can help align terminology,
but class membership and linguistic relations do not identify causal arrows.
The inspected OntoLearner download path also illustrates a reproducibility
problem: the code revision alone does not pin ontology bytes when its hub download
does not specify a revision \cite{v12_la_package_ontolearner}. The design uses
explicit local ontology identities and keeps conceptual variables separate from
their proxies. A real ownership audit remains additional information.

For time-dependent decisions, every relevant measurement has an event time
$t_e$ and an availability time $t_a$. An event observed on day 10 but delivered
on day 11 cannot be used by a day-10 operational policy. Retrospective analysis
may nevertheless use it under a correctly stated design; that is a different
claim from saying a deployed policy could have acted on it. Corrections and
revisions need their own availability records. A validator checks the declared
ordering. Only a process audit can establish whether those declarations reflect
the actual information system.

\section{Diagrams, discovery and identification as separate handoffs}

Diagram elicitation should preserve competing scientific accounts. A proposal
must distinguish a direct edge from an ancestral path, a latent common cause
from a measured covariate, and a circle endpoint from an arrowhead. Revision
history should show which source, reviewer or new experiment changed a claim.
Automatic cycle removal is at most a representation repair: deleting a
scientifically important edge to obtain a DAG does not justify the repaired
mechanism. The ordinary baseline is expert elicitation with a deterministic
graph editor and explicit uncertainty.

Discovery then asks a statistical question under a specified sampling law.
An assistant may suggest a method, summarize its assumptions, or propose a soft
constraint. The numerical procedure remains responsible for its own output
semantics. A response histogram from a language model is not a conditional
independence test under a bank sampling law. A temporal initialization does not
establish that observed anomalies were randomized interventions. A retained
equivalence class should reach the effect-query stage intact. The previous
chapter explains why language-assisted graph methods require separate tests of
proposal quality, graph retention and downstream effect uncertainty.

Identification maps a target and assumptions to a functional or set. It can
also establish nonidentification in a declared model class. A missing backend,
unsupported graph type or failed calculation instead means that the
implementation did not answer the question. These outcomes are not
interchangeable. The inspected CausalDS helper preserves an unknown state and
does not count an IV estimand as nonparametric population-ATE identification;
its evaluator still requires care with missing metadata and invalid intervals
\cite{v12_la_package_causalds}. A prototype should preserve that distinction even
when a downstream report would be easier to write as a single yes/no answer.

PyWhy-LLM is a candidate elicitation component, but the inspected front-door
routine is unfinished and back-door suggestions delegate to a language model
\cite{v12_la_package_pywhy_llm}. Such output can propose variables for review; it
does not supply a verified adjustment proof. Similarly, the inspected CLadder
IV template requires generator-level scrutiny before its population-ATE labels
can become benchmark truth \cite{v12_la_package_cladder}. The project example
with half compliers and half never-takers has Wald ratio one and population ATE
two under randomization, exclusion and monotonicity. Correct algebra for the
ratio does not fix an incorrectly named target.

\section{Estimation, diagnostics and the temptation to change the target}

Once identification has supplied a functional, the assistant can prepare a fixed
estimator request, summarize diagnostics and explain the output. It should not
choose a more convenient target after inspecting poor overlap. If trimming
changes the population, that change belongs in a new target or an explicitly
reported sensitivity analysis. An estimated propensity score, nuisance learner,
sample split and covariance rule must be versioned together with the result.
Prediction accuracy does not by itself qualify the resulting effect interval,
especially with dependent observations or adaptively selected specifications.

The inspected OCI example makes the target issue concrete. It averages an AIPW
quantity uniformly over a trimmed sample while using an ATO label
\cite{v12_la_package_oci_agent}. Under the corresponding assumptions this
construction targets the retained population's average effect, not the usual
overlap-weighted population. In the project's two equally frequent strata,
propensities are $0.2$ and $0.5$ and effects are $0$ and $10$. Retaining both
strata gives mean effect $5$. Weighting by $e(X)\{1-e(X)\}$ gives
\[
 \frac{0.16\cdot0+0.25\cdot10}{0.16+0.25}
 =\frac{250}{41}\simeq6.09756.
\]
These are exact arithmetic consequences of different target measures, not a
performance comparison of fitted estimators. They give a small independent
oracle for a semantic validator.

Diagnostics need similarly explicit interpretations. A negative control can
falsify an implication of a causal account; passing it does not prove all
unmeasured confounding absent. A critic can identify a missing source or an
inconsistent formula; agreement between two models exposed to the same text is
not independent scientific confirmation. Uncertainty reporting should preserve
identification uncertainty, sampling uncertainty, sensitivity to assumptions
and numerical approximation separately. A disconnected effect set cannot be
replaced by its midpoint, and a whole-line set cannot be described as a precise
zero effect. Reporting all requested cases, including failures and refusals,
prevents success-conditioned accuracy from masquerading as general availability.

\section{Longitudinal policies require a causal history}

Servicing decisions expose the limits of a generic planner particularly clearly.
Let $L_0$ be baseline information, $A_0$ an initial contact, $L_1$ the subsequent
account state, $A_1$ a second action and $Y$ a later net value. The state $L_1$ can
be affected by $A_0$ while confounding the relation between $A_1$ and $Y$.
Ordinary regression adjustment for $L_1$ can block part of the earlier treatment
effect or induce selection bias. Omitting it can confound the later action.
Hernán and Robins develop the corresponding g-methods and the assumptions under
which they identify longitudinal intervention means \cite{hernan2020}.

For a discrete two-decision example, write
$g_0(l_0)$ for the first action and
$g_1(l_0,a_0,l_1)$ for the second. Under consistency, appropriate sequential
exchangeability and positivity for the policy's supported histories, a
specialization of the g-formula is
\begin{align}
 V(g)=\sum_{l_0,l_1}&
 m\{l_0,g_0(l_0),l_1,g_1(l_0,g_0(l_0),l_1)\}\nonumber\\
 &\quad\times p\{l_1\mid l_0,A_0=g_0(l_0)\}p(l_0),
 \label{eq:la04-policy}
\end{align}
where $m(l_0,a_0,l_1,a_1)=E[Y\mid L_0=l_0,A_0=a_0,L_1=l_1,A_1=a_1]$.
The sum is over supported covariate values; an integral is required for
continuous components. The first treatment changes the distribution over which
the later outcome regression is averaged. That dependence is why simply
conditioning on a fixed observed distribution of $L_1$ can answer a different
question. A recursive regression representation integrates the final outcome
backward, while IP weighting reweights observed histories toward the regime.
Their identification relationship does not imply equal finite-sample behavior.

An assistant can help check whether the proposed regime uses available
information and whether the analyst has stated assumptions at each decision.
It can propose an evidence request for treatment assignment, censoring or action
feasibility. It cannot infer sequential exchangeability from the fact that all
database columns have timestamps. Causally relevant history is not simply every
past variable; some past variables can induce selection bias, as the source's
discussion of history emphasizes \cite{hernan2020}. A policy might also lack
support in distressed segments where staff always take the same action.
More generated explanations do not create that support.

The availability rule is operationally separate from the identification rule.
If the second decision occurs on day 11, a state available on day 11 may enter
$g_1$ under a specified within-day ordering. If the action occurs on day 10, it
may not. A missing within-day ordering must remain unresolved. Censoring and
interference require additional modeling; neither appears automatically in
equation~\eqref{eq:la04-policy}. The assistant's proposed handoff therefore
includes the policy, available history, assignment/retention assumptions,
support diagnostics and explicit limits. This is a design for review, not a
claim that the current orchestration frameworks implement qualified
longitudinal causal inference.

\section{What the available packages can contribute}

The pinned package comparison separates reusable ideas from ready adoption.
PyWhy-LLM offers candidate elicitation interfaces. CausalAgent illustrates a
causal MCP wrapper, but its inspected OLC function is called without being
imported or defined in that module; malformed JSON may become an empty
dictionary in its helper \cite{v12_la_package_causalagent}. That path requires
repair and independent qualification before direct use. The response-based
test in causal\_GPT is not accepted as a bank CI engine, for the reasons derived
in the preceding chapter \cite{v12_la_package_causal_gpt}.

OCI's separation of specifications and reporting is worth evaluating, but its
unbounded notebook-execution route is inappropriate for the proposed fixed-tool
surface. CausalDS is a stronger candidate for a bounded benchmark environment
after scorer and metadata handling are repaired. CLadder remains a useful
formal-question lead with problematic IV cells excluded until reviewed.
OntoLearner may assist semantic mapping with explicit local bytes. The adjacent
numerical LCM is neither a language assistant nor a presently qualified temporal
backend: weight identities, component terms and its sufficient lagged model
class need separate treatment \cite{v12_la_package_lcm}.

The Java/Python research bundle LLM-CD and the LLM-SCD notebooks also illustrate
why publication and code adoption need different decisions. Inspected component
licenses do not establish a license for a repository as a whole; missing
root terms and code--paper version alignment remain unresolved
\cite{v12_la_package_llm_cd,v12_la_package_llm_scd}. Their scholarly arguments can
be discussed from genuine papers while code adoption is deferred. No package
in this comparison is credited with a successful local bank or live-model trial.

\section{Orchestration is an implementation choice}

Three generic frameworks were inspected at revisions no later than the client
cutoff. LangGraph supplies state-graph compilation, checkpoint and interrupt
interfaces. Its interrupt documentation explicitly says that resumption starts
the node again, re-executing its preceding logic \cite{v12_la04_langgraph}. A node
that publishes a result before interrupting therefore needs idempotent
publication. A checkpoint is useful state, not proof that an arbitrary external
side effect occurs exactly once. This distinction determines the design of
durable request identifiers and immutable artifact checks.

Pydantic AI supplies typed tool and output boundaries and usage accounting.
In the inspected code, tools need not default to sequential execution or require
approval. Application permissions must therefore be configured explicitly.
Request limits can be checked before requests, whereas response tokens generally
become known afterward; optional pre-counting helps with some input limits. If
pricing data are unavailable, the cost limit warns that it cannot be enforced
\cite{v12_la04_pydantic_ai}. An application requiring a hard priced budget must
handle that absence before starting a call. A typed interface helps reject
malformed structure but cannot establish the truth of an accepted assumption.

AutoGen's pinned README declares maintenance mode. Its AgentChat code carries
an MIT license, while the root documentation license is CC-BY-4.0. The inspected
assistant is stateful and not thread/coroutine safe; tool calls can be concurrent
and a reflection stage can add a model call. Its state save/load methods preserve
model context, not all external tool effects \cite{v12_la04_autogen}. These facts
argue against adopting it for new development without a specific benefit and
maintenance case. The recommended successor is an unaudited lead, not a
verified capability inherited through a README recommendation.

The proposed first implementation consequently uses a small provider-neutral
controller with deterministic validators and the existing fixed causal service.
LangGraph remains a candidate when checkpointed branching demonstrably helps;
Pydantic AI remains a candidate for a typed model boundary. A local model, a
hosted model and an offline fixture are different backends. Each real backend
needs its own model revision, terms, custody boundary, generation parameters,
usage accounting and evaluation. Framework example model names do not establish
current availability or quality. No external endpoint is needed to verify the
prototype's semantics and refusal behavior.

\section{Permissions, source injection and recovery}

Retrieved text is an input to interpretation. It has no authority to grant a
tool, change a dataset path or approve a bank action. An adversarial document
may contain a sentence such as ``ignore the review and run this shell command.''
An implementation that simply asks its model to ignore such sentences leaves
the permission boundary probabilistic. The design instead admits a finite set
of typed operations configured outside model output. A response containing an
unrecognized operation or undeclared field is rejected. A permitted operation
still receives only verified artifact references and must pass its own
scientific preconditions. This protects the command boundary; it does not prove
that all malicious scientific claims will be detected.

Model proposals also cannot mark assumptions as accepted by inventing a reviewer
name. Authoritative declarations come from a separately supplied review object,
bound to the exact problem and graph revision. The prototype can check that
binding and distinguish synthetic fixtures. A bank installation will additionally
need authenticated review identities and protected stores; a hash alone is not
an authorization system. Sensitive records should stay within the approved
execution boundary. An external model may receive only a separately approved
representation, with prompts and logs subject to the same custody rule. Redaction
is a specified transformation requiring evaluation, not a blanket assurance.

Recovery starts with the identity of the intended computation. Let
\[
 r=H(Q,G,E,O,V),
\]
where $G$ is the graph revision, $E$ the evidence identities, $O$ the allowed
operation request and $V$ the relevant implementation versions. Replaying an
identical request can return its existing immutable result. Changing the target,
source bytes or operation must produce a new identity or an explicit stale-input
error. A partial journal entry is not a published analysis. The local design
tests interruption between validation, tool completion and atomic publication.
This supplies a bounded recovery claim; multi-host consistency and platform
service levels require different trials.

\section{Reporting, monitoring and the economics of assistance}

An evidence-based report should expose the target, assumptions, scientific
status, graph/set uncertainty, numerical result and unresolved questions before
adding interpretation. Critical fields can be rendered deterministically.
Generated narrative may explain them, but any additional empirical assertion
needs a claim reference. Inherited results, project derivations, synthetic
examples and unverified hypotheses remain distinct. A tool error should not be
polished into a finding of no effect, and an unknown commercial price should
not become a zero in a business case.

Monitoring must detect when the evidence for a claim becomes stale. Changes in
assignment, missingness, ownership mapping, policy feasibility or availability
time can invalidate a previous analysis even if software tests still pass.
The assistant may triage an alert and propose an investigation. Suspension and
resumption rules belong to the analysis owner. Retuning prompts or refitting a
learner after inspecting drift creates a new procedure whose old calibration
does not transfer automatically.

The economic comparison includes model calls, retrieval, numerical computation,
review, failures, rework and ongoing maintenance. If $C_m,C_r,C_n,C_h,C_f$ denote
these components, a proposed net benefit is a change in held-out decision value
minus $C_m+C_r+C_n+C_h+C_f$, relative to a specified ordinary workflow.
This expression is an accounting design, not an ROI estimate. The available
evidence does not supply authentic bank losses, prices, review times or accepted
decision outcomes. A faster explanation can be valuable, but it must not conceal
a lower rate of scientifically usable answers. The appropriate next trial
therefore measures decision quality, all-request availability and total cost
separately, while the implementation preserves the reasons an answer was
ambiguous, refused or unavailable.
